\section{Antes de empezar}
\subsection*{Cómo entras}
\begin{enumerate}
  \item Entra a Traza FI con tu correo institucional (\texttt{@alumno.uaemex.mx}).
  \item En el menú, abre \emph{Titulación} y luego \emph{Mi trámite de licenciatura}.
  \item En \emph{¿Cómo quieres titularte?} abre \emph{Sin trabajo escrito} o \emph{Con trabajo escrito} y elige tu opción. Antes de iniciar puedes leer su página y descargar este manual; también puedes comparar las opciones en \emph{Compara las opciones: descarga sus manuales}.
\end{enumerate}
\begin{tencuenta}
\begin{itemize}
  \item El avance se mide por \textbf{etapas} («Etapa 2 de 8»); dentro de cada etapa, la tarjeta \emph{Tu siguiente paso} te dice qué hacer ahora.
  \item Cada documento va en PDF escaneado, legible y a color (no fotografías); el tamaño máximo lo indica la pantalla.
  \item Solo puede haber \textbf{una modalidad activa}. Si cambias de opción, antes se te explica qué pasará con la que tienes y se te pide confirmarlo.
  \item Los formatos se llenan con tus datos y se firman electrónicamente; no necesitas imprimirlos.
\end{itemize}
\end{tencuenta}
\begin{algosalemal}
\begin{itemize}
  \item \textbf{Te devolvieron un documento.} Arriba de tu trámite aparece \emph{Falta corregir un documento} con la observación. Corrige el PDF y vuelve a subirlo desde el mismo renglón.
  \item \textbf{Tu asesor no encuentra el correo.} Desde tu trámite puedes reenviarle el enlace; revisa que su correo esté bien escrito.
  \item \textbf{Elegiste la opción equivocada.} Al final de tu trámite abre \emph{¿Elegiste otra modalidad o quieres dejar este trámite?}. Si nadie más ha intervenido, el trámite \textbf{se borra}; si ya intervino alguien (asesor, comité, Titulación), \textbf{queda en pausa} en \emph{Abandonadas} y conservas todo.
  \item \textbf{Quieres volver a un trámite en pausa.} En \emph{Mi trámite de licenciatura}, sección \emph{Abandonadas}, pulsa \boton{Retomar}: sigue justo donde se quedó.
  \item Si una pantalla muestra un error que no entiendes, toma nota del folio de tu trámite y escribe a Apoyo a Titulación.
\end{itemize}
\end{algosalemal}
